\begin{textsample}{POS Dim 1 – vem\_ed\_05 – Score 6.00 – vem\_ed\_05/t019.txt}  \label{ex:f1_pos_104}
Intercâmbios Virtuais: pesquisa 2020 Ao término dos Projetos Colaborativos Internacionais realizados no segundo semestre de 2020, estudantes e professores responsáveis pelos Intercâmbios Virtuais nas 20 Fatecs participantes \textbf{foram} convidados a responder a uma pesquisa de percepção.
Dos 34 professores, 25 responderam; dos 833 alunos, aproximadamente 400 retornaram o questionário.
A \textbf{seguir}, uma síntese dos principais resultados.
ALUNOS 79 \% indicariam a experiência para outros alunos 84 \% consideram \textbf{que} a \textbf{competência} no \textbf{idioma} estrangeiro melhorou 72 \% avaliam \textbf{interação} com colegas estrangeiros como"ótima"ou"boa"86 \% avaliam \textbf{interação} com colegas brasileiros"ótima"ou"boa"63 \% acham \textbf{que} o Intercâmbio virtual melhora desempenho acadêmico PROFESSORES 100 \% recomendariam a experiência para outros docentes 88 \% afirmam \textbf{que} a \textbf{competência} no \textbf{idioma} estrangeiro melhorou. 88 \% avaliam como"ótima"e 12 \% como"boa"a \textbf{interação} com colegas estrangeiros. 92 \% \textbf{observaram} \textbf{que} a \textbf{interação} dos alunos melhorou com a inserção dos projetos \textbf{colaborativos} nas atividades de aula 72 \% concordam totalmente \textbf{que} as colaborações contribuem para pesquisas na área do professor 92 \% concordam totalmente \textbf{que} o \textbf{intercâmbio} virtual possibilita novas atividades acadêmicas a \textbf{partir} da experiência do parceiro internacional.
Ferramentas digitais mais \textbf{utilizadas} Whatsapp E-mail Zoom

% matched lemmas: colaborativo, competência, idioma, interação, intercâmbio, observar, partir, que, seguir, ser, utilizar
\end{textsample}
